\SourceLeaf{270}%
\noindent folglich mit Benutzung von 10):

\[
\tag*{(F)} df(x) = f(dx)\, g(x).
\]

16. Aus der geometrischen Bedeutung von $f(x)$ geht hervor, dass $f(0) = 0$ ist. Mithin darf man

\[
\tag*{12)} f(dx) = \alpha\, dx
\]

setzen, wo $\alpha$ eine Constante bedeutet. Wird zur Abkürzung

\[
f(x) = y
\]

\Needspace{6\baselineskip}\noindent gesetzt, so ist nach (F)

\[
\tag*{13)} \frac{dy}{dx} = \alpha\, g(x),
\]

also

\[
\tag*{14)} g(x) = \frac{1}{\alpha}\frac{dy}{dx}.
\]

Aus 13) folgt durch Differentiation:

\[
\frac{d^2y}{dx^2} = \alpha\, \frac{dg(x)}{dx},
\]

mithin nach (E), 12) und 14):

\[
\frac{d^2y}{dx^2} = -\alpha^2\, \frac{1 - \dfrac{1}{\alpha^2}\dfrac{dy^2}{dx^2}}{y},
\]

oder.

\[
\tag*{15)} y\, \frac{d^2y}{dx^2} - \left(\frac{dy}{dx}\right)^2 + \alpha^2 = 0.
\]

Um diese Differentialgleichung zu integriren, multiplicire man sie mit

\[
\frac{2}{y^3}\, \frac{dy}{dx},
\]

wodurch sie in

\[
\frac{2\dfrac{dy}{dx}\dfrac{d^2y}{dx^2}}{y^2} - \frac{2\left(\dfrac{dy}{dx}\right)^3}{y^3} + 2\alpha^2\frac{\dfrac{dy}{dx}}{y^3} = 0
\]

übergeht. Mithin ist

\[
\frac{1}{y^2}\left(\frac{dy}{dx}\right)^2 - \frac{\alpha^2}{y^2} = \text{Const.}
\]

Wird die Constante mit $-k^2$ bezeichnet, so kommt

\[
\left(\frac{dy}{dx}\right)^2 - \alpha^2 + k^2 y^2 = 0,
\]

\Needspace{9\baselineskip}\noindent woraus

\[
dx = \frac{dy}{\sqrt{\alpha^2 - k^2 y^2}} = \frac{\dfrac{k}{\alpha}\, dy}{k\sqrt{1 - \dfrac{k^2}{\alpha^2} y^2}}
\]

und

\[
kx = \arcsin \frac{k}{\alpha} y + \text{const.}
\]
